\section{Discussion}
\label{sec:discussion}

\para{Answer to the research question.}
A secret language without a shared key is possible (\thmref{thm:keyless}; \secref{sec:e2}), but it is not a \emph{language} in the emergent sense. It is cryptography that both agents must already know how to run (\propref{prop:learn}). Anything the agents can work out tacitly from shared knowledge, an equally informed monitor can work out too (\corref{cor:schelling}). In \eone, the partner and a same-model auditor were statistically indistinguishable ($p=0.75$), while the weaker \qwensmall monitor lost by 20 discordant pairs to 0. Tacit codes threaten only monitors that are weaker or less informed than the agents.

\para{Where the secret goes.}
Removing the key moves secrecy to the agents' \emph{private randomness}, which they exercise through \emph{interaction}. It moves the coordination burden to \emph{prior common knowledge of the protocol}. This knowledge does not have to be secret. Published protocols such as those of \citet{vaikuntanathan2026undetectable} and \citet{pang2026codetta} can provide it, and as capable models read such papers, the protocol becomes a Kerckhoffs focal point that is cheap for colluders and visible to monitors at the same time. \citet{rippin2026tool} suggest that agents with tools can already implement such systems.

\para{Implications for oversight.}
Our results give concrete guidance. First, monitors should run the same weights with the same context as the agents. This defeats every tacit and weight-specific channel (\secref{sec:e1}; Corollaries~\ref{cor:schelling} and~\ref{cor:weights}). Second, passive transcript auditing cannot catch channels built on cryptographic key exchange, even in principle. Third, such channels are fragile against active measures: light resampling or paraphrasing (\secref{sec:warden}), impersonation probes (\propref{prop:mitm}), and having the monitor supply the agents' sampling randomness. If the monitor knows the sampling coins, \lemref{lem:sim} applies to every message and no secrecy is possible.

\para{Comparison with prior work.}
Our \eone results agree with \citet{jarviniemi2025subversion} and \citet{zolkowski2025early}: without a pre-agreed scheme, LLMs fall back on near-plaintext hiding that monitors catch (AUROC 0.93). Our \etwo cost of about 7 tokens per key bit for a 0.5B model is in the same order of magnitude as Codetta's roughly 80k tokens for a 128-bit key. The comparison is not like-for-like, because Codetta includes error correction and a more conservative design. Prior analyses of active wardens model them as binary symmetric channels~\citep{roger2023preventing,perry2025robust}. Our desynchronization finding shows that this model can underestimate the damage to entropy-gated channels by more than an order of magnitude.

\para{Surprises.}
The Kerckhoffs auditor beat the partner (32\% vs.\ 21\%). The ``extract'' framing of the auditor prompt seems to matter more than the role itself. The active-warden collapse came from desynchronization, not bit flips, which we did not anticipate when planning the experiment.

\para{Limitations.}
We used no frontier or API models. \eone uses open models up to 14B parameters, which are weaker steganographers than frontier models. \corref{cor:schelling} is a theorem and does not depend on scale, but the absolute coordination rates do. \etwo is toy cryptography: the main setting uses a 256-bit group (one 2048-bit run checks the cost), has no error correction, and operates on token IDs rather than text. We checked undetectability only against the tests we ran (NLL, rank, monobit, ASCII, Legendre). The honest null set has only 80 messages, so the AUROC CIs are about $\pm 0.07$. The \publicplain condition also showed a marginal rank/NLL shift ($p \approx 0.04$--$0.06$, not significant after correction). Our warden resamples tokens independently given the original prefix. Real paraphrasers rewrite at the meaning level, which is even more destructive to token-level channels, and we do not test meaning-level channels~\citep{zhou2026conceptual}. In \eone, the scheme taxonomy is coarse and LLM-classified, the partner and auditor prompts differ in framing as well as role, each condition has 100 payloads and a single sender model, and the shared-scheme sender followed the scheme only 36\% of the time. Finally, the new propositions are proved informally (\appref{app:proofs}); \propref{prop:learn} is stated for a computationally bounded receiver and has not been machine-checked.

\para{Broader impact.}
This work describes how covert channels between AI agents can be built and broken. The constructions we implement are already public~\citep{vonahn2004public,vaikuntanathan2026undetectable,pang2026codetta}. Our contribution is mainly defensive: it identifies which monitor capabilities defeat which channels, and it adds a cheap active defense.
